\section{Finite-rank metric constraints and a semisimple control}
\label{met-section}

Additional vacua can compensate the curvature of an arithmetic ground-state
line without introducing negative norms.  They do not, however, automatically
preserve its physical state family or its arithmetic observable.  We establish
two finite-rank restrictions and then give a positive rank-two solution of the
metric equations that illustrates the distinction.

\subsection{The metric equation and its determinant}
\label{met-conventions}

Let $z=\sigma+i\eta$ be an ordinary chiral parameter.  In an intrinsic
holomorphic frame, write the positive Hermitian vacuum metric as $g$ and set
\begin{equation}
 A_z=g^{-1}\partial_zg,\qquad
 \mathcal K=\partial_{\bar z}(g^{-1}\partial_zg),\qquad
 C^{\dagger_g}=g^{-1}C^*g.
 \label{met-connection}
\end{equation}
Here $C^*$ denotes conjugate transpose.  Our convention for the ordinary
chiral $tt^*$ equation is
\begin{equation}
 \partial_{\bar z}C=0,\qquad
 \mathcal K=[C,C^{\dagger_g}].
 \label{met-ttstar}
\end{equation}
Indeed $[D_z,D_{\bar z}]=-\mathcal K$ for
$D_z=\partial_z+A_z$, $D_{\bar z}=\partial_{\bar z}$, and
\eqref{met-ttstar} gives flatness of
$D_z+\zeta^{-1}C$ and $D_{\bar z}+\zeta C^{\dagger_g}$.
The vacuum-metric interpretation originates in topological--antitopological
fusion~\cite{CecottiVafa}; the parameter-multiplet hypothesis is essential to
the supersymmetric Berry-connection derivation~\cite{SonnerTong}.
The existence of four formal supercharges alone does not identify a specified
parameter family with \eqref{met-ttstar}.

For a finite set of primes $S$, the coherent-state benchmark has
\begin{equation}
 h(\sigma)=\Lambda_S(\sigma),\qquad
 \Lambda_S(z)=\pi^{-z/2}\Gamma(z/2)
       \prod_{p\in S}(1-p^{-z})^{-1},\qquad \re z>0.
 \label{met-arithmetic-metric}
\end{equation}
Writing $u=\log h$ gives
\begin{equation}
 V(\sigma):=u''(\sigma)
 =\frac14\psi_1(\sigma/2)
  +\sum_{p\in S}\frac{(\log p)^2p^{-\sigma}}
                         {(1-p^{-\sigma})^2}>0,
 \qquad
 \partial_z\partial_{\bar z}\log h(\re z)=\frac V4.
 \label{met-variance}
\end{equation}
This physical norm is distinct from the squared holomorphic amplitude
$|\Lambda_S(z)|^2$.

\begin{proposition}[Determinant constraint]
\label{met-determinant-proposition}
For a finite-rank positive solution of \eqref{met-ttstar},
$\partial_z\partial_{\bar z}\log\det g=0$.
Consequently $\det g=|d(z)|^2$ locally for a nonvanishing holomorphic $d$.
In particular, the determinant cannot equal
$h(\re z)^a|d(z)|^2$ for a nonzero real $a$.
\end{proposition}
\begin{proof}
The finite-dimensional trace gives
\[
 \partial_z\partial_{\bar z}\log\det g
 =\operatorname{tr}\mathcal K
 =\operatorname{tr}[C,C^{\dagger_g}]=0.
\]
A real harmonic function is locally the real part of a holomorphic function,
which proves the factorization.  The proposed determinant instead has
logarithmic curvature $aV/4\ne0$ by \eqref{met-variance}.
\end{proof}

Thus tensoring a rank-$N$ solution with the original arithmetic line creates
determinant curvature $NV/4$.  Prescribing $h$ as one component is a different
operation: $g=\operatorname{diag}(h,h^{-1})$ is positive and has determinant
one.  Determinant cancellation is a necessary condition, not the full
differential compatibility condition.

\subsection{Ambient holomorphicity is a stronger restriction}
\label{met-ambient}

\begin{proposition}[Finite ambient-holomorphic obstruction]
\label{met-ambient-proposition}
Let $\mathcal H$ be a fixed positive Hilbert space and let
$T(z):\C^N\to\mathcal H$ be an injective Hilbert-norm holomorphic frame,
with $N<\infty$.  Set $g=T^*T$ and
$P=Tg^{-1}T^*$, the orthogonal projector onto its span.
If this induced metric satisfies \eqref{met-ttstar}, then $P$ is constant on
each connected parameter domain and $\mathcal K=0$.
\end{proposition}
\begin{proof}
With $T'=\partial_zT$, holomorphicity gives
$\partial_zg=T^*T'$ and $\partial_{\bar z}g=T'^*T$.  Hence
\begin{align}
 \mathcal K
 &=-g^{-1}T'^*Tg^{-1}T^*T'+g^{-1}T'^*T'
   =g^{-1}T'^*(1-P)T',
 \label{met-second-fundamental}\\
 \operatorname{tr}\mathcal K
 &=\norm{(1-P)T'g^{-1/2}}_{\mathrm{HS}}^2.
 \label{met-positive-trace}
\end{align}
The trace must also vanish by Proposition~\ref{met-determinant-proposition}.
Thus $(1-P)T'=0$ and $\mathcal K=0$.  Direct differentiation yields
$\partial_zP=(1-P)T'g^{-1}T^*=0$; taking the adjoint also gives
$\partial_{\bar z}P=0$.
\end{proof}

This excludes a direct finite-dimensional holomorphic enlargement of the
complete Morse--prime coherent family.  In its fixed positive Hilbert space,
that family is
\begin{equation}
 \Psi_z=e^{zX/2}e^{-e^y/2},\qquad
 X=\frac y2-\sum_{p\in S}n_p\log p-\frac{\log\pi}{2},
 \qquad n_p\in\{0,1,2,\ldots\}.
 \label{met-coherent-family}
\end{equation}
Projection to all $n_p=0$ sends states at distinct positive real parameters
$\sigma_j$ to nonzero scalar multiples of
$e^{\sigma_jy/4}e^{-e^y/2}$.  These are linearly independent: division by
$e^{-e^y/2}$ and successive differentiation at one point gives a Vandermonde
system.  Arbitrarily large independent sets exist on every open real
interval.  Consequently no constant finite-dimensional subspace can contain
this family, as Proposition~\ref{met-ambient-proposition} would require.

The assumption $\partial_{\bar z}T=0$ must not be substituted for the weaker
physical condition $P\partial_{\bar z}T=0$.  An intrinsic holomorphic vacuum
frame may have nonzero antiholomorphic derivatives in the massive complement.
Such derivatives are allowed in physical $tt^*$ geometry and invalidate the
positive-trace calculation \eqref{met-positive-trace}.  Infinite rank likewise
requires separate trace and domain arguments.

The corresponding line-bundle identity makes the remaining freedom explicit.
For an intrinsic holomorphic section $v$, let
$h_v=\ip{v}{v}_g$ and $\beta=(1-P_L)D_zv$, where $P_L$ projects onto its line.
In a holomorphic normal frame, differentiating $\log h_v$ gives
\begin{equation}
 \partial_z\partial_{\bar z}\log h_v
 =\frac{\ip{v}{\mathcal Kv}_g+\norm{\beta}_g^2}{h_v}
 =\frac{\norm{C^{\dagger_g}v}_g^2-\norm{Cv}_g^2
                  +\norm{\beta}_g^2}{h_v}.
 \label{met-gauss-line}
\end{equation}
Thus positive line curvature can be supplied by chiral mixing and variation
inside a larger vacuum bundle.  It is not itself a rank-two obstruction.

\subsection{A local positive semisimple metric}
\label{met-semisimple}

Rank-two positivity and reality conditions are studied in~\cite{Takahashi}.
The following explicit construction verifies the metric equations directly;
it makes no existence assertion about a microscopic quantum theory.

\begin{proposition}[Real-axis metric control]
\label{met-control-proposition}
There is a neighborhood of $z=1/2$ carrying a positive rank-two solution of
\eqref{met-ttstar}, with semisimple $C_z$, whose first diagonal component on
the real axis is $\Lambda_S(\sigma)$.
\end{proposition}
\begin{proof}
For $r=\re w>0$, set
\begin{equation}
 g_*(w,\bar w)=
 \begin{pmatrix}\coth(2r)&0\\0&\tanh(2r)\end{pmatrix},
 \qquad
 C_* =\begin{pmatrix}0&1\\1&0\end{pmatrix}.
 \label{met-coth-control}
\end{equation}
The first curvature entry equals the first commutator entry because
\begin{equation}
 \frac14\frac{d^2}{dr^2}\log\coth(2r)
 =\frac{4\cosh(4r)}{\sinh^2(4r)}
 =\coth^2(2r)-\tanh^2(2r).
 \label{met-coth-verification}
\end{equation}
The second entries have opposite sign.  The metric is positive with determinant
one, while $C_*$ has distinct eigenvalues $\pm1$.  For the constant bilinear
pairing $\eta_0=\left(\begin{smallmatrix}0&1\\1&0\end{smallmatrix}\right)$,
$\eta_0C_*$ is symmetric and
$(\eta_0^{-1}g_*)\overline{(\eta_0^{-1}g_*)}=1$.

At $\sigma=1/2$, $h>1$ and $h'<0$.  Indeed
$\Gamma(1/4)>\int_0^1t^{-3/4}(1-t)\,dt=16/5$,
$\pi^{1/4}<2$, and each Euler factor exceeds one.  Also
$\psi(1/4)=-\gamma-\pi/2-3\log2<0$, while the logarithmic derivatives of
the Euler factors are negative.  A local branch therefore exists for
\begin{equation}
 w(z)=\frac14\log\frac{\Lambda_S(z)+1}{\Lambda_S(z)-1},
 \qquad
 g_{\mathrm{new}}(z,\bar z)=g_*(w(z),\overline{w(z)}),
 \qquad C_z=w'(z)C_*.
 \label{met-pullback}
\end{equation}
After shrinking the neighborhood, $\re w>0$ and $w'\ne0$ there.  Holomorphic
pullback multiplies both sides of \eqref{met-ttstar} by $|w'|^2$.
On the real axis $2w=\operatorname{arccoth}h$, so
$(g_{\mathrm{new}})_{11}(\sigma,\sigma)=\coth(2w(\sigma))=h(\sigma)$.
The eigenvalues of $C_z$ are the distinct values $\pm w'(z)$.
\end{proof}

The complex extension has changed:
$(g_{\mathrm{new}})_{11}(z,\bar z)$ is generally not $h(\re z)$.
Only a real-axis diagonal norm has been retained.  The construction does not
specify the physical states \eqref{met-coherent-family}, their complete overlap
$\Lambda_S((\bar z+w)/2)$, a self-adjoint arithmetic observable $X$, or a
Hamiltonian with the proposed metric as its vacuum metric.  Semisimplicity
and the displayed algebraic reality condition also do not establish the
additional ultraviolet, infrared, and boundary data of a massive quantum
field theory.  The domain is local, with no all-prime or global half-plane
continuation asserted.

\begin{remark}[Keeping the original complex extension]
\label{met-companion-remark}
For $g=\operatorname{diag}(h(\sigma),h(\sigma)^{-1})$ and
$\kappa=V/4$, the companion choices
$C=\left(\begin{smallmatrix}0&1\\q&0\end{smallmatrix}\right)$ and
$C=\left(\begin{smallmatrix}0&q\\1&0\end{smallmatrix}\right)$ require,
respectively,
\[
 |q(z)|^2=R_+(\sigma):=h^4-\kappa h^2,
 \qquad
 |q(z)|^2=R_-(\sigma):=\kappa h^{-2}+h^{-4}.
\]
Holomorphic $q$ forces $RR''-(R')^2=0$ wherever it is nonzero.  Stirling's
formula gives $u''=1/(2\sigma)+O(\sigma^{-2})$ and superexponential growth of
$h$, whence $(\log R_+)''=2/\sigma+O(\sigma^{-2})$ and
$(\log R_-)''=-1/\sigma+O(\sigma^{-2})$ at infinity.
Real analyticity on $(0,\infty)$ excludes either identity even on an open
patch; identically zero $q$ is excluded by the same large-$\sigma$ behavior.
This is a test of the stated frame and companion normalization, not every
rank-two Higgs field.  The changed extension in \eqref{met-pullback} avoids
these hypotheses.
\end{remark}

The control is deliberately permissive: any positive real-analytic norm
greater than one with nonzero derivative can be inserted into the same local
pullback.  In particular, local metric compatibility does not select the
arithmetic source normalization or its first normal logarithmic derivative.
The next physical construction must preserve an explicitly defined complete
pairing, rather than only one metric component, and derive its normalization
from the theory.  None of the metric results above supplies the Weil pole
pairing or proves positivity of the complete Weil form.
